% =====================================================================
%  7. ANÁLISIS CRÍTICO DEL PROCESO DE AUDITORÍA
%  Responsable: INTEGRANTE 4
%  ESTADO: REDACTADO (1 de agosto de 2026) — ~2.800 palabras
%
%  REGLA APLICADA: toda afirmación crítica va sostenida por una fuente.
%  Criticar sin evidencia resta tanto como no criticar. Y donde el
%  proceso hizo bien las cosas, se dice.
%
%  El conflicto de evidencia (7.3) NO se resuelve artificialmente. Si
%  alguien lo reescribe tomando partido, se pierde lo que más vale del
%  capítulo.
% =====================================================================

\section{Análisis crítico del proceso de auditoría}
\label{sec:critico}

Las secciones anteriores describieron el proceso de auditoría y sistematizaron sus
resultados. Corresponde ahora evaluarlo: determinar en qué medida el proceso
reunió las condiciones que la práctica auditora exige y qué consecuencias tuvo
sobre la conducta efectiva de la organización examinada.

El criterio que guía esta sección es doble. Toda afirmación crítica se sostiene en
una fuente identificable, porque un juicio sin evidencia carece de valor en un
informe de auditoría. Y donde el proceso resolvió correctamente, se consigna con
la misma claridad con que se consignan sus deficiencias: un análisis que solo
encuentra defectos no es riguroso, es tendencioso.

% ---------------------------------------------------------------------
\subsection{El problema de la independencia}

La objeción de mayor calado contra la evaluación analizada se refiere a su
financiamiento. La organización Rights Action calificó el proceso de fundamental e
irrevocablemente viciado e inaceptable, precisamente por haber sido costeado por
la empresa auditada \parencite{rightsactionsfcritica}.

La objeción no puede desestimarse. La ISO 19011:2018 establece la independencia
como principio rector de la auditoría y exige que el auditor esté libre de sesgo y
de conflicto de intereses \parencite{iso2018iso19011}. Cuando quien paga es
también quien resulta examinado, existe un conflicto de intereses estructural, con
independencia de la integridad personal de los profesionales implicados.

Frente a ello debe ponderarse, sin embargo, el diseño institucional efectivamente
adoptado, expuesto en la sección \ref{sec:auditoria}. El encargo no procedió de la
dirección de la empresa sino de un Comité Directivo independiente, que fue además
el órgano que seleccionó a la firma auditora. El Memorando de Entendimiento fijó
por escrito los principios de transparencia, independencia e inclusividad. Y el
informe se publicó íntegro, en versión bilingüe, incluidos los hallazgos
desfavorables \parencite{oncommonground2010hria}. Ese último punto no es menor: si
el financiamiento hubiera determinado el resultado, el documento no contendría la
constatación de que la operación se autorizó sin consulta previa adecuada, que es
el reproche más grave que puede formularse a un proyecto extractivo en territorio
indígena.

Conviene por tanto distinguir dos nociones que la crítica de Rights Action tiende
a fundir: la \textbf{independencia financiera} y la \textbf{independencia de
criterio}. La primera no se dio; la segunda parece haberse preservado en lo
sustantivo, a juzgar por el contenido del informe. Ahora bien, la distinción tiene
un límite, y es donde la objeción recupera su fuerza: la independencia de criterio
opera sobre lo que el auditor decide escribir, pero no necesariamente sobre lo
que decide investigar, ni sobre el alcance que acepta, ni sobre el mandato de
seguimiento que se le reconoce. Y es precisamente en esos tres planos —alcance,
profundidad y seguimiento— donde la sección \ref{sec:hallazgos} localizó las
mayores debilidades del proceso.

La pregunta pertinente no es, en consecuencia, si la evaluación estaba viciada,
sino qué modelo alternativo habría sido preferible. Las opciones realmente
existentes son cuatro. La \textbf{fiscalización estatal} aporta potestad
sancionadora, pero depende de la capacidad institucional del Estado, que en este
caso no había activado ningún procedimiento de consulta. La \textbf{supervisión
multilateral}, del tipo ejercido por la Compliance Advisor Ombudsman, es
independiente del operador, pero su alcance se limita al cumplimiento de las
políticas del organismo financiador \parencite{cao2006marlin}. El \textbf{panel
independiente financiado por terceros} —modelo del panel internacional promovido
desde el ámbito académico que encargó el estudio de salud— elimina la dependencia
económica, pero carece de acceso garantizado a las instalaciones y a la
documentación interna de la empresa. Y la \textbf{certificación de tercera parte
bajo estándar auditable}, como el sistema IRMA, combina auditoría independiente con
requisitos verificables y umbrales cuantificados, aunque su contratación también
la costea la empresa \parencite{irmasfstandard}.

Ninguno de los cuatro modelos elimina por completo la tensión. Lo que diferencia a
unos de otros es el conjunto de salvaguardas y, sobre todo, la existencia de un
mecanismo de verificación posterior independiente. Ese es el elemento que faltó
aquí, y es el que se propone incorporar en la sección \ref{sec:plan}.

% ---------------------------------------------------------------------
\subsection{Limitaciones metodológicas declaradas}

El propio equipo auditor reconoció que la participación del municipio de Sipacapa
y de los sectores opositores a la mina resultó insuficiente, y calificó en
consecuencia sus hallazgos sobre impactos como parciales
\parencite{oncommonground2010hria}.

Debe empezarse por el mérito, porque lo tiene. Declarar las propias limitaciones
es una buena práctica de auditoría y responde al principio de presentación
imparcial de la ISO 19011:2018, que obliga a informar con veracidad y exactitud
incluidos los obstáculos encontrados \parencite{iso2018iso19011}. Un informe que
hubiera ocultado esta carencia sería peor, no mejor.

Dicho esto, el peso de la limitación es considerable y conviene dimensionarlo con
precisión. No se trata de una laguna cuantitativa —haber entrevistado a menos
personas de las deseables— sino de un \textbf{sesgo de selección} orientado en la
dirección que más importa. En una evaluación cuyo objeto es determinar si una
operación afecta derechos, la población que sostiene haber sido afectada
constituye la fuente primaria de evidencia. Su ausencia no reduce
proporcionalmente la fiabilidad del resultado: la compromete de manera asimétrica,
porque los hallazgos sobre impactos negativos quedan sistemáticamente
subrepresentados mientras los relativos al desempeño formal de la empresa
conservan su respaldo documental.

Enlaza esto con el principio de enfoque basado en la evidencia, según el cual las
conclusiones deben derivarse de un muestreo apropiado. Un muestreo del que se
excluye involuntariamente al grupo más crítico no es apropiado en sentido técnico,
por muy declarada que esté la exclusión.

A ello se añade la segunda limitación: la revisión por pares a cargo de
International Alert, prevista en el diseño del proceso, no llegó a ejecutarse
según lo planeado. La revisión por pares es el mecanismo mediante el cual un
trabajo técnico se somete al escrutinio de terceros competentes antes de su
publicación. Su ausencia significa que el informe se publicó sin haber sido
contrastado por nadie ajeno al equipo que lo produjo. En un proceso ya
cuestionado por su dependencia financiera, la pérdida de la única salvaguarda
técnica externa prevista resulta particularmente inoportuna. El expediente no
aclara las causas de esa omisión.

% ---------------------------------------------------------------------
\subsection{El conflicto de evidencia sobre la contaminación}

La sección \ref{sec:hallazgos} estableció que la evidencia sobre afectación
efectiva de la calidad del agua es contradictoria. Corresponde ahora examinar por
qué lo es, sin resolver artificialmente la contradicción.

De un lado, \textcite{basu2010scitotenv} hallaron concentraciones
significativamente más altas de mercurio, cobre, arsénico y zinc en orina en
quienes residían más cerca de la mina que en el grupo de comparación, y
\textcite{etech2010water} concluyeron que los efectos adversos sobre el ambiente
podían haber comenzado ya. Del otro, la empresa y determinados monitoreos
oficiales sostuvieron que no existía contaminación atribuible a la operación. Y
el propio estudio epidemiológico matizó su hallazgo: la mayoría de los metales
se detectaron por debajo de los valores asociados a daño clínico, y sus autores
concluyeron que hacían falta estudios más robustos.

Ninguna de las tres posiciones es descartable con la información pública
disponible, y la razón es metodológica. Los estudios difieren en la
\textbf{matriz analizada} —biomonitoreo humano frente a agua superficial—, en los
\textbf{puntos de muestreo}, en los \textbf{períodos} y en los
\textbf{laboratorios} empleados. Dos mediciones que no comparten protocolo no son
comparables, aunque ambas sean correctas en su propio marco.

Existe además una limitación de fondo, señalada ya en la sección
\ref{sec:hallazgos}: un estudio de biomonitoreo acredita exposición, no fuente.
La presencia de arsénico en orina no identifica su procedencia, y en un territorio
mineralizado el aporte geogénico natural puede ser considerable. Establecer
atribución exige comparar contra una línea base previa al inicio de la operación.

Y aquí reside el punto crítico, ya documentado en la sección
\ref{sec:hallazgos}: la línea base disponible cubrió apenas ocho o nueve meses,
se limitó a dos manantiales en materia de agua subterránea y no estableció las
direcciones de flujo \parencite{etech2010water}. Sin ella, la controversia carece de
método de resolución. La consecuencia es que el desacuerdo sobre la contaminación
no es un problema de honestidad de las partes ni de calidad de los laboratorios:
es una \textbf{consecuencia previsible del diseño del sistema de monitoreo}. Se
construyó una operación de gran escala sobre un territorio habitado sin dotarse
del instrumento que permitiría, años después, demostrar la propia inocuidad.

Esta es la razón por la que la sección \ref{sec:hallazgos} clasifica la
insuficiencia de la línea base como no conformidad mayor y la afectación efectiva
como observación. El fallo verificable no está en el resultado, está en el
dispositivo. Lo que haría converger la evidencia —línea base pública, monitoreo
participativo con cadena de custodia, laboratorio acreditado y validación por
terceros— se propone en la sección \ref{sec:plan}.

% ---------------------------------------------------------------------
\subsection{Consecuencias institucionales y eficacia real del proceso}

La evaluación no operó en el vacío. Se inserta en una cadena de pronunciamientos
de órganos con distinto grado de autoridad, cuyo recorrido conviene reconstruir
para valorar su eficacia agregada.

En \textbf{2005-2006}, tras la denuncia de la organización Madre Selva, la
Compliance Advisor Ombudsman concluyó que la Corporación Financiera Internacional
no había aplicado adecuadamente sus políticas de salvaguarda social y ambiental al
aprobar el préstamo al proyecto, y recomendó un proceso participativo de diálogo.
El caso se cerró en mayo de 2006 \parencite{cao2006marlin}.

En \textbf{2010} se publicó la evaluación de derechos humanos, y ese mismo año la
Comisión de Expertos de la Organización Internacional del Trabajo instó a
suspender actividades hasta dar cumplimiento a la consulta del Convenio 169
\parencite{oit1989convenio169}.

El \textbf{20 de mayo de 2010} la Comisión Interamericana de Derechos Humanos
otorgó medidas cautelares a favor de miembros de dieciocho comunidades del pueblo
maya de San Marcos, solicitando la suspensión de la mina. En
\textbf{diciembre de 2011} modificó esas medidas y retiró la solicitud de
suspensión, reorientándolas hacia el acceso a agua apta para consumo humano
\parencite{cidh2010medidas}.

En \textbf{2017} la operación cerró por agotamiento de reservas, con un plan de
recuperación ambiental que incluyó destrucción de cianuro y relleno del tajo
\parencite{goldcorp2023marlin}.

De esta secuencia se desprende la valoración más incómoda de este informe. Durante
el período en que estuvieron vigentes unas medidas cautelares internacionales que
solicitaban expresamente la suspensión de la actividad, la actividad no se
suspendió: la propia empresa hizo constar que las operaciones continuaron con
normalidad a lo largo de todo el procedimiento ante la Comisión. La afirmación es
verídica y no fue desmentida.

Cabe entonces formular la pregunta directamente: ¿modificó la auditoría la
conducta de la organización auditada? La respuesta honesta es que la modificó de
forma parcial y en los márgenes. Hubo efectos verificables —la planta de
tratamiento de agua, el monitoreo comunitario, la transición a relaves filtrados en
2013, la integración de aprendizajes a las políticas corporativas globales— y
todos ellos merecen registrarse. Pero la no conformidad mayor del expediente, la
ausencia de consulta previa, nunca se cerró; la recomendación de detener la
expansión hasta consultar debidamente no se cumplió; y la operación concluyó su
ciclo íntegro conforme a lo previsto.

Ello sugiere una conclusión que trasciende el caso: una auditoría voluntaria,
carente de potestad sancionadora y sin mecanismo de verificación independiente del
seguimiento, produce mejoras incrementales en aquellos ámbitos donde la mejora es
compatible con la continuidad de la operación, pero no altera las decisiones
estructurales del proyecto. No es un instrumento inútil —los cambios documentados
son reales— pero tampoco es un sustituto de la regulación. Confundir ambas cosas
es el riesgo principal que este tipo de procesos comporta, y probablemente el
fundamento más sólido de la crítica de Rights Action, aun cuando su formulación
resulte excesiva.

% ---------------------------------------------------------------------
\subsection{Contraste con el régimen peruano}

Procede finalmente preguntarse qué habría ocurrido si los mismos hechos se
hubieran producido bajo el ordenamiento peruano expuesto en la sección
\ref{sec:marco}. El ejercicio es hipotético, pero resulta ilustrativo de la
diferencia entre auditoría y fiscalización.

En primer lugar, la \textbf{consulta previa} dispone en el Perú de un
procedimiento reglado, con etapas y plazos definidos por la Ley 29785, que
desarrolla internamente la obligación del Convenio 169
\parencite{congreso2011ley29785}. La ausencia de consulta no habría sido una
omisión sin cauce, sino el incumplimiento de un trámite exigible, susceptible de
impugnación en sede administrativa y judicial.

En segundo lugar, el hallazgo relativo al monitoreo ambiental habría quedado bajo
competencia del Organismo de Evaluación y Fiscalización Ambiental, que ejerce
funciones de supervisión, fiscalización, control y sanción sobre la mediana y
gran minería \parencite{congreso2009ley29325}. La diferencia con lo efectivamente
ocurrido es de naturaleza: donde la evaluación de derechos humanos formuló una
recomendación, el organismo fiscalizador habría podido dictar una medida
correctiva de cumplimiento obligatorio, con el respaldo del artículo 39 del
Reglamento de Protección y Gestión Ambiental Minera
\parencite{minem2014ds040}, e imponer sanciones ante el incumplimiento.

En tercer lugar, la certificación ambiental previa habría correspondido al Servicio
Nacional de Certificación Ambiental para las Inversiones Sostenibles, con las
exigencias de línea base del Sistema Nacional de Evaluación de Impacto Ambiental
\parencite{congreso2001ley27446}. La ausencia de línea base pública —el fallo
sistémico identificado en la sección \ref{sec:hallazgos}— habría sido un defecto
del expediente de certificación, no una carencia advertida cinco años después.

Sería sin embargo un error concluir que el modelo peruano resuelve la
conflictividad socioambiental minera. El caso de Espinar demuestra lo contrario:
allí existían un aparato fiscalizador, una mesa de diálogo y un monitoreo sanitario
ambiental participativo de considerable escala técnica —313 puntos de monitoreo y
40 sustancias analizadas por punto— y aun así el conflicto persistió
\parencite{minam2013espinar}. Las experiencias de Las Bambas y Tía María apuntan
en el mismo sentido.

La diferencia entre uno y otro modelo es, por tanto, de \textbf{exigibilidad
jurídica}, no de resultados garantizados. El régimen peruano ofrece cauces
formales y potestad sancionadora de los que el caso guatemalteco careció; lo que
no ofrece es la garantía de que esos cauces se activen con oportunidad, ni de que
las comunidades afectadas los perciban como legítimos. Esa distinción —entre
disponer del instrumento y usarlo eficazmente— es la que conviene retener, y es la
que fundamenta que la propuesta de la sección \ref{sec:plan} no se limite a
invocar el cumplimiento normativo, sino que incorpore mecanismos de verificación
independiente y de participación comunitaria efectiva.
